% page 592 du fac-simile (image p-591 du scan)
% bloc 154.9 x 246.3 mm ; marge gauche 8.0 mm
\pgc{-12.700mm}{0.000mm}{
\l{\cel{2}584}
\l{}
\l{\sou{tomato}. (\sou{bot.}) I. Planto ek la familio \textquotedbl{}solanei\textquotedbl{}, kun frukti}
\l{\cel{3}brilante-reda. L.\cel{1}\sou{solanum lycopersicum}. - II. La frukto}
\l{\cel{3}di ta planto, uzata por facar dishi o salado-forme.-DEFRS}
\l{}
\l{\sou{tombo}. I. Foso superkovrita de quaza tablo de petro, de mar-}
\l{\cel{3}moro, qua kontenas mortinto. - II. Monumento funerala}
\l{\cel{3}super la foso qua kontenas mortinto. -\cel{2}EFS.}
\l{}
\l{\sou{tombako}.\cel{2}Aloyuro de 90\% latuno e 10\% zinko. - DEFIRS.}
\l{}
\l{\sou{tono}. (\sou{psik}.) Flexo di la voco, qua indikas la diversa senti}
\l{\cel{3}di la psiko. - II. (\sou{muziko)}\cel{2}(a) Intervalo inter du noti,}
\l{\cel{3}(b) Singla de la gami en qui muzikajo esas kompozebla,}
\l{\cel{3}indikata per la noto komencala qua determinas la repartis-}
\l{\cel{3}eso di la intervali e mi-intervali inter du noti. - III.}
\l{\cel{3}(\sou{fotogr., pikt.}) Grado di la sama koloro segun ke lu esas}
\l{\cel{3}plu o min obskura, plu o min briloza. - IV. (\sou{fiziol}.)}
\l{\cel{3}Kontrakturo mikra e permananta di ula muskuli, precipue}
\l{\cel{3}olti di la sfinteri. - DEFIRS.}
\l{}
\l{\sou{tondar}. (\sou{trans}.)Tranchar reze (la lano, la pili). - FI.}
\l{}
\l{\sou{tondrar}. (\sou{netrans}.) (\sou{meteorol.}) (\sou{aludante fulmino-stroko})}
\l{\cel{3}Produktar bruiso resonanta, son-eganta. - DEFI.}
\l{}
\l{\sou{tonika}. I. (\sou{gram.})\cel{1}\sou{Acento tonika}\cel{1}: Plualtigo (che la Anti-}
\l{\cel{3}qui) o plu-intensigo (che la Moderni) di la sono lor la}
\l{\cel{3}prononco di silabo determinita. - II. (\sou{muziko}).\cel{1}\sou{Noto}\cel{1}to-}
\l{\cel{3}\sou{nika}\cel{1}: Noto qua karakterizas la tono, la gamo ye qua ario,}
\l{\cel{3}melodio, muzikajo kompozesis - ta de qua lu komencas e}
\l{\cel{3}qua donas a lu sua nomo. -\cel{2}DEFIS.}
\l{}
\l{\sou{tonoplasto}. (\sou{biol.}) Parto permananta e diferenciita di la}
\l{\cel{3}protoplasmo, qua formacas la parieto di la vakuoli, e pre-}
\l{\cel{3}cipue di la vakuoli kontraktebla. - DESIF.}
\l{}
\l{\sou{tonsilo}. (\sou{anat.}) Korpo glandatra, mandelo-forma, situita an}
\l{\cel{3}singla latero di la istmo di la guturo, faringo, laringo.}
\l{\cel{3}- DEFIS.}
\l{}
\l{\sou{tonsilito.}\cel{2}(\sou{patol.}) Inflamuro di la tonsili. -}
\l{}
\l{\sou{tonsuro}. (\sou{religio katolika}). I. Krono quan onu facas sur la}
\l{\cel{3}kapo di la kleriki, diakoni, sacerdoti, monaki, e c.}
\l{\cel{3}per razar parto di la hari. - II. Ceremonio dum qua la}
\l{\cel{3}episkopo grantas ad ulu la grado unesma di la klerikeso,}
\l{\cel{3}per razar la hari de-sur la somito di la kapo. - DEFIRS.}
\l{}
\l{\sou{tontino}. Asociitaro de personi qui kolokas kapitalo komuna,}
\l{\cel{3}por obtenar de lu, ye epoko determinita, rento dum-viva,}
\l{\cel{3}qua reversionesos (pos singla morto) a la transvivinti,}
\l{\cel{3}di qui la porciono kreskas proporcione. - DEFIRS.}
\l{}
\l{topo.\cel{2}(\sou{aludante navo}).\cel{2}Platformo muntita salie cirkum, ed}
\l{\cel{3}an la extremajo supra di basa masto qua trairas lu. E.}
}
